% =====================================================================
%  Scott A. Wolpert,
%  "Lectures and notes: Mirzakhani's volume recursion and approach for the
%   Witten-Kontsevich theorem on moduli tautological intersection numbers"
%  arXiv:1108.0174v1 [math.DG] (31 Jul 2011) —— 中文译注本
%
%  编译：LuaLaTeX（本机 xdvipdfmx 不稳定，统一用 lualatex，编译 3 遍）
%  用法：powershell -ExecutionPolicy Bypass -File build.ps1
%
%  本底本的特殊性（译者必读）：
%   * 原文的**编号公式只有 4 条**：(1)(2)(3) 在 p.18，(4) 在 p.21；正文引用了它们（"formula (1)"、
%     "formulas (3)"）。这 4 条必须用 plain TeX 的 \eqno{(N)} 手工复刻编号，写法：
%         \[ H(x,y)=... \eqno{(1)} \]
%     两行共用一个编号时用 \[ \begin{aligned} ... \end{aligned} \eqno{(2)} \]。
%     其余公式**一律不编号**，禁止 \tag 与带编号的 equation/align。
%   * 原文的编号对象是**定理类陈述**：Theorem 1--15（其中 12 号是 Definition），
%     编号在全书连续，必须手工复刻（\begin{thm}{N}）。这是本项目的验收锚点。
%   * 图 1--11 为原文矢量图 500 dpi 截取（0 位图），图注为译者所译。
%   * % [Tutor Note: ...] 为译者批注，不进入排版，汇总见 译注汇总.md。
% =====================================================================
\documentclass[11pt,a4paper]{article}

\usepackage[UTF8]{ctex}
\usepackage{amsmath,amssymb,amsthm,mathtools}
\usepackage{bm}
\usepackage{graphicx}
\usepackage[a4paper,left=2.5cm,right=2.5cm,top=2.6cm,bottom=2.7cm]{geometry}
\usepackage[dvipsnames]{xcolor}
\usepackage{enumitem}
\usepackage{booktabs}
\usepackage{array}
\usepackage{multirow}
\usepackage{longtable}
\usepackage{multicol}
\usepackage{fancyhdr}
\usepackage{caption}
\usepackage{tcolorbox}
\usepackage{tikz}
\usepackage[hidelinks,unicode]{hyperref}
\hypersetup{pdftitle={讲义与注记：Mirzakhani 的体积递推与 Witten–Kontsevich 定理（中译注本）},
  pdfauthor={Scott A. Wolpert; 中译/批注: 本译本}}

\captionsetup{font=small,labelsep=quad}
\setlength{\emergencystretch}{3em}
\allowdisplaybreaks[1]
\graphicspath{{figures/}{../figures/}}

% ---------- 页眉 ----------
\pagestyle{fancy}
\fancyhf{}
\fancyhead[L]{\small\itshape Mirzakhani's volume recursion \& Witten-Kontsevich}
\fancyhead[R]{\small\thepage}
\renewcommand{\headrulewidth}{0.4pt}

% ---------- 定理类环境：编号一律由原文手工给出（原文 1--15 连续编号） ----------
% 用法：\begin{thm}{3} ... \end{thm}  /  \begin{prop}{11} / \begin{defn}{12} / \begin{cor}{N}
\newenvironment{thm}[1]{\par\medskip\noindent\textbf{定理 #1.}\ }{\par\medskip}
\newenvironment{prop}[1]{\par\medskip\noindent\textbf{命题 #1.}\ }{\par\medskip}
\newenvironment{defn}[1]{\par\medskip\noindent\textbf{定义 #1.}\ }{\par\medskip}
\newenvironment{corl}[1]{\par\medskip\noindent\textbf{推论 #1.}\ }{\par\medskip}
\newenvironment{lemmaenv}[1]{\par\medskip\noindent\textbf{引理 #1.}\ }{\par\medskip}

% 原文中不编号的 "Theorem." 段（如原文 p.3 的两条）——只有粗体标签、无编号
\newenvironment{thmplain}{\par\medskip\noindent\textbf{定理.}\ }{\par\medskip}
\newenvironment{thmplaincite}[1]{\par\medskip\noindent\textbf{定理.}\ [#1].\ }{\par\medskip}

% 原文 Proof. 环境
\renewcommand{\proofname}{证明}

% 原文的粗体/斜体行内小标题（如 "The papers." "First paper brief."）——
% 中文译为 \para{...}，独立成段、粗体、不缩进。
\newcommand{\para}[1]{\par\medskip\noindent\textbf{#1}\quad\ignorespaces}

% 原文图标题（手工编号，不用 \caption 自动编号）
\newcommand{\figtitle}[2]{\caption*{\small\textbf{图 #1：}#2}}
% 原文表标题（本项目底本 0 表，保留宏备用）
\newcommand{\tabtitle}[2]{\caption*{\small\textbf{表 #1：}#2}}

% 参考文献条目（Lead 手工转录，保留英文原貌）
\newcommand{\refitem}[2]{\par\hangindent=1.8em\hangafter=1\noindent [#1]~#2}

% 原文页码定位（批注/校对用；不进入排版）
\newcommand{\origp}[1]{}

% ---------- 常用记号（只能用这里定义过的宏，不要自造 \newcommand） ----------
\newcommand{\R}{\mathbb{R}}
\newcommand{\C}{\mathbb{C}}
\newcommand{\Z}{\mathbb{Z}}
\newcommand{\N}{\mathbb{N}}
\newcommand{\Q}{\mathbb{Q}}
\newcommand{\PP}{\mathbb{P}}
% 花体：模空间 M、Teichmüller T、曲线 C 等
\newcommand{\cA}{\mathcal{A}}
\newcommand{\cB}{\mathcal{B}}
\newcommand{\cC}{\mathcal{C}}
\newcommand{\cD}{\mathcal{D}}
\newcommand{\cE}{\mathcal{E}}
\newcommand{\cF}{\mathcal{F}}
\newcommand{\cG}{\mathcal{G}}
\newcommand{\cH}{\mathcal{H}}
\newcommand{\cI}{\mathcal{I}}
\newcommand{\cJ}{\mathcal{J}}
\newcommand{\cK}{\mathcal{K}}
\newcommand{\cL}{\mathcal{L}}
\newcommand{\cM}{\mathcal{M}}
\newcommand{\cN}{\mathcal{N}}
\newcommand{\cO}{\mathcal{O}}
\newcommand{\cP}{\mathcal{P}}
\newcommand{\cQ}{\mathcal{Q}}
\newcommand{\cR}{\mathcal{R}}
\newcommand{\cS}{\mathcal{S}}
\newcommand{\cT}{\mathcal{T}}
\newcommand{\cU}{\mathcal{U}}
\newcommand{\cV}{\mathcal{V}}
\newcommand{\cW}{\mathcal{W}}
\newcommand{\cX}{\mathcal{X}}
\newcommand{\cY}{\mathcal{Y}}
\newcommand{\cZ}{\mathcal{Z}}
% 本论文专用记号的简写
\newcommand{\T}{\mathcal{T}}                     % Teichmüller space（译者常用裸 \T）
\newcommand{\M}{\mathcal{M}}                     % moduli space（译者常用裸 \M）
\newcommand{\Tg}{\mathcal{T}_g}                 % Teichmüller space
\newcommand{\Tgn}{\mathcal{T}_{g,n}}
\newcommand{\Mg}{\mathcal{M}_g}                 % moduli space
\newcommand{\Mgn}{\mathcal{M}_{g,n}}
\newcommand{\cMgn}{\mathcal{M}_{g,n}}
\newcommand{\dimc}{\dim_{\C}}
\newcommand{\Mbar}{\overline{\mathcal{M}}}      % Deligne-Mumford 紧化
\newcommand{\Mgb}{\overline{\mathcal{M}}_g}
\newcommand{\Mgnb}{\overline{\mathcal{M}}_{g,n}}
\newcommand{\Tbar}{\overline{\mathcal{T}}}      % augmented Teichmüller space
\newcommand{\Tgnb}{\overline{\mathcal{T}}_{g,n}}
\newcommand{\MCG}{\mathrm{MCG}}                 % mapping class group
\newcommand{\WP}{\mathrm{WP}}                   % Weil-Petersson
\newcommand{\Vol}{\operatorname{Vol}}
\newcommand{\vol}{\operatorname{vol}}
\newcommand{\CV}{\operatorname{CV}}             % covolume
\newcommand{\Int}{\operatorname{Int}}
\newcommand{\Area}{\operatorname{Area}}
\newcommand{\len}{\ell}
\newcommand{\dl}{\mathrm{d}\ell}
\newcommand{\dtau}{\mathrm{d}\tau}
\newcommand{\dd}{\mathrm{d}}
\newcommand{\ee}{\mathrm{e}}
\newcommand{\ii}{\mathrm{i}}
\newcommand{\Tr}{\operatorname{Tr}}
\newcommand{\tr}{\operatorname{tr}}
\newcommand{\diag}{\operatorname{diag}}
\newcommand{\sgn}{\operatorname{sgn}}
\newcommand{\rk}{\operatorname{rk}}
\newcommand{\Res}{\operatorname{Res}}
\newcommand{\Id}{\mathrm{Id}}
\newcommand{\Stab}{\mathrm{Stab}}
\newcommand{\ie}{\textit{i.e.}}
\newcommand{\eg}{\textit{e.g.}}
\newcommand{\cf}{\textit{cf.}}
\newcommand{\half}{\tfrac{1}{2}}
\newcommand{\wt}[1]{\widetilde{#1}}
\newcommand{\what}[1]{\widehat{#1}}
